\section{Conclusion}
\label{sec:conclusion}

This paper presented the Variance-Aware Adaptive Bit-Width Quantisation (VABQ) framework for energy-efficient air quality monitoring on edge IoT devices. By jointly optimising sensor-level ADC resolution and inference-level model precision through a unified variance-entropy criterion, VABQ achieves a 64.6\% reduction in total system energy consumption while maintaining pollutant classification accuracy above 96.1\% and signal reconstruction NRMSE below 2.4\%. Experiments on the DALTON dataset---spanning 28 indoor sites, 42 occupants, and three months of continuous multi-channel data at 1~Hz---demonstrate the practical viability of the approach on ESP32-based hardware with minimal memory and computational overhead.

The dual-layer design, which bridges the traditionally separate domains of sensor data compression and neural network quantisation, constitutes the principal novelty and offers a template for energy-aware sensing-inference co-optimisation across diverse IoT monitoring applications. The deterministic, threshold-based sensor-level controller combined with standard post-training quantisation ensures reproducibility and ease of deployment---significant advantages over more complex hardware-aware NAS approaches.

Future work will extend the framework to outdoor air quality networks with heterogeneous sensor modalities, integrate adaptive window sizing through change-point detection, and explore end-to-end quantisation-aware training across both sensing and inference layers. Additionally, validation across diverse geographical regions and climate zones will strengthen the generalisability claims of the VABQ methodology.
